\section*{Overview}

Bipartite matching is the cleanest example of turning a structural graph question into an algorithmic primitive. Many
``pair these objects'' problems are really asking for the largest set of disjoint edges between a left side and a right
side.

\section*{When to Use It}

Use bipartite matching when:

\begin{itemize}
  \item the objects naturally split into two groups,
  \item each allowed pairing links one left item to one right item,
  \item every item can be used only a limited number of times, often exactly once.
\end{itemize}

Examples include job assignment, student-project pairing, rook placement, or scheduling with compatibility edges.

\section*{Core Idea}

A matching is a set of edges with no shared endpoints. The algorithm grows the matching by finding augmenting paths:
paths that alternate between unmatched and matched edges and start and end at free vertices.

Flipping the status of the edges on such a path increases the matching size by one.

\section*{Key Insight}

The alternating-path view explains both correctness and implementation. Hopcroft-Karp speeds things up by finding many
shortest augmenting paths in one BFS/DFS phase instead of increasing the matching by only one edge at a time.

\section*{Operations / Main Technique}

\begin{itemize}
  \item build the bipartite graph from left to right,
  \item BFS from free left vertices to compute distance layers,
  \item DFS through those layers to find augmenting paths,
  \item repeat until no free right vertex is reachable.
\end{itemize}

\paragraph{Worked example.}
If \(L_1\) is free, \(R_3\) is matched to \(L_4\), and there is a path
\[
L_1 \to R_3 \to L_4 \to R_2,
\]
where \(R_2\) is free, flipping matched and unmatched edges along that path increases the total matching size by one.

\section*{Correctness Intuition}

An augmenting path always increases the matching by one because it starts and ends at free vertices and alternates
edge status in between. When no augmenting path exists, the current matching is maximum. Hopcroft-Karp is still based
on that fact; it only batches shortest augmenting paths together.

\section*{Complexity Analysis}

Hopcroft-Karp runs in \(O(E \sqrt{V})\), which is usually much better than repeated DFS matching on larger graphs.

\section*{Implementation}

The code uses Hopcroft-Karp with:

\begin{itemize}
  \item adjacency lists from the left side,
  \item distance labels for the BFS phase,
  \item two match arrays, one per side.
\end{itemize}

\section*{Common Pitfalls}

\begin{itemize}
  \item Forgetting that the graph must be bipartite before this model applies.
  \item Mixing left-side and right-side indexing in the match arrays.
  \item Reaching for general max flow when plain bipartite matching is enough.
  \item Using Kuhn's algorithm on input sizes that really need Hopcroft-Karp.
\end{itemize}

\section*{Variants / Extensions}

\begin{itemize}
  \item Minimum vertex cover in bipartite graphs via Konig's theorem.
  \item Weighted assignment, which leads to Hungarian or min-cost flow instead.
  \item Capacity on one side by splitting or by a small flow reduction.
\end{itemize}

\section*{Practice Problems}

\begin{itemize}
  \item Maximum bipartite matching on compatibility edges.
  \item Place the maximum number of non-attacking objects on a constrained grid.
  \item Recover a minimum vertex cover after finding a maximum matching.
\end{itemize}

\section*{References}

\begin{itemize}
  \item \href{https://cp-algorithms.com/graph/kuhn_maximum_bipartite_matching.html}{cp-algorithms: Bipartite Matching}.
  \item \href{https://usaco.guide/adv/bipartite-matching}{USACO Guide: Bipartite Matching}.
  \item \href{https://codeforces.com/catalog}{Codeforces Catalog}.
\end{itemize}
